% !TEX root = ../main.tex
% Ejemplo de conversacion de tres turnos (IT-106).
%
% El caso esta elegido para que se vea el mecanismo y no para que quede bonito:
% la titulacion NO aparece en ninguna de las dos preguntas del usuario, solo en
% la respuesta del asistente. Es el unico ejemplo en el que se distingue mirar
% las preguntas de mirar tambien la respuesta, que es la correccion de IT-106.
\begin{figure}[H]
  \centering
  \resizebox{\textwidth}{!}{%
    \begin{tikzpicture}[
      turno/.style={draw=diagNodoBorde, fill=diagNodo, align=left,
                    font=\small, text width=6.6cm, inner sep=6pt,
                    rounded corners=3pt},
      asistente/.style={turno, fill=white},
      sujeto/.style={draw=diagGrupoBorde, fill=diagGrupo, align=center,
                     font=\small, text width=5.8cm, inner sep=5pt,
                     rounded corners=3pt},
      rotulo/.style={font=\small\itshape, anchor=west},
      hilo/.style={-{Latex[length=2mm]}, semithick, draw=diagGrupoBorde!70!black},
    ]
      \node[rotulo] at (0,0.9) {Conversación};
      \node[rotulo] at (8.4,0.9) {De qué se está hablando};

      % --- Turno 1 ---
      \node[turno, anchor=north west] (p1) at (0,0.4)
        {\textbf{Estudiante:} me gustan los videojuegos};
      \node[asistente, anchor=north west] (r1) at (0,-1.0)
        {\textbf{Asistente:} te encaja el \textbf{Grado en Ingeniería
         Informática}, que tiene asignaturas de programación y gráficos};
      \node[sujeto, anchor=north west] (s1) at (8.4,-1.0)
        {Grado en Ingeniería Informática\\{\footnotesize sale de la
         \textbf{respuesta}}};

      % --- Turno 2 ---
      \node[turno, anchor=north west] (p2) at (0,-3.4)
        {\textbf{Estudiante:} ¿y qué asignaturas tiene en primero?};
      \node[sujeto, anchor=north west] (s2) at (8.4,-3.4)
        {Grado en Ingeniería Informática\\{\footnotesize se \textbf{hereda}: la
         pregunta no nombra ninguna}};

      % --- Turno 3 ---
      \node[turno, anchor=north west] (p3) at (0,-5.6)
        {\textbf{Estudiante:} ¿y en Mecánica?};
      \node[sujeto, anchor=north west] (s3) at (8.4,-5.6)
        {Grado en Ingeniería Mecánica\\{\footnotesize \textbf{cambia}: lo que la
         pregunta nombra manda}};

      \draw[hilo] (s1.south) -- (s2.north);
      \draw[hilo] (s2.south) -- (s3.north);
    \end{tikzpicture}%
  }
    \caption{Resolución y evolución del sujeto conversacional a lo largo de tres turnos de diálogo.}
  \label{fig:conversacion_tres_turnos}
\end{figure}
